\documentclass{article}
% Shared layout. Keep benchmark prose and result-specific plots outside this folder.
\PassOptionsToPackage{numbers}{natbib}
\usepackage[preprint]{neurips_2026}
\usepackage[T1]{fontenc}
\usepackage[hidelinks]{hyperref}
\usepackage{booktabs,amsmath,microtype,pgfplots}
\pgfplotsset{compat=1.17}
\usetikzlibrary{patterns}
\definecolor{okblue}{HTML}{0072B2}
\definecolor{okorange}{HTML}{D55E00}
\pgfplotsset{every axis/.append style={font=\small, tick label style={font=\small}, width=\linewidth, height=5cm, grid=major}}

\newcommand{\PaperTitle}{Benchmark Skateboard: Signed Rotations and Flight-Time Constraints}
\newcommand{\PaperAbstract}{Benchmark Skateboard tests consistency between signed rotations, projectile arithmetic, and landing constraints under an explicit simplified model. Its 72 public tasks cross 24 trick definitions with three gravity settings. We evaluate scripted controls through the benchmark runner and conduct a simulated population experiment with numeric and rotation corruption. A glossary-assisted analytic control reaches full credit on every task. Reversing rotation signs retains a mean score of 0.780, showing why partial credit must accompany complete-answer accuracy. Numeric perturbations reveal a sharp tolerance boundary with an absolute rounding-floor exception. Simulated repair layers separate arithmetic recovery from consistency with corrupted rotations. All results are generated locally without paid inference. These findings characterize the measurement instrument; they do not establish learned-model skill or real-world skateboarding ability.}
\newcommand{\TaskCount}{72}
\newcommand{\TrickCount}{24}
\newcommand{\SimAgentCount}{20}
\newcommand{\SimMultiplier}{1.1}
\newcommand{\BootstrapDraws}{2000}
\newcommand{\ConditionCount}{313}
\newcommand{\ObservationCount}{22536}
\newcommand{\AuthorOrcid}{0009-0005-9964-6841}
\newcommand{\OracleMean}{1.000 [1.000, 1.000]}
\newcommand{\OracleFull}{1.000 [0.949, 1.000]}
\newcommand{\ConstantMean}{0.020 [0.000, 0.050]}
\newcommand{\ConstantFull}{0.000 [0.000, 0.051]}
\newcommand{\AlwaysMean}{0.805 [0.760, 0.847]}
\newcommand{\AlwaysFull}{0.250 [0.164, 0.361]}
\newcommand{\ReverseMean}{0.780 [0.760, 0.810]}
\newcommand{\ReverseFull}{0.083 [0.039, 0.170]}
\newcommand{\Simnone}{0.401 [0.384, 0.416]}
\newcommand{\Simphysics}{0.880 [0.873, 0.887]}
\newcommand{\Gainphysics}{0.480 [0.463, 0.498]}
\newcommand{\Simboolean}{0.379 [0.364, 0.393]}
\newcommand{\Gainboolean}{-0.021 [-0.025, -0.018]}
\newcommand{\SensitivityOutside}{0.243 [0.240, 0.248]}

\title{\PaperTitle}
\author{Pooria Arab\\Independent researcher\\ORCID: \AuthorOrcid}
\begin{document}
\maketitle
\begin{abstract}\PaperAbstract\end{abstract}
\section{Why this matters}
A correct physical answer must connect several kinds of information. A solver must interpret a direction, calculate motion, and apply a condition to the result. Success on arithmetic alone does not establish this connection. A fluent explanation can also hide a reversed axis or a mistaken landing judgment. Benchmark Skateboard makes these dependencies explicit in a small, auditable setting.
The construct is \emph{consistency with stated rotation and flight constraints}. A task describes a skateboard trick, a coordinate convention, an upward launch speed, gravity, and a minimum flight duration. The answer reports signed turns, airtime, peak height, final board orientation, and whether the duration is sufficient. The equations and decision rules appear in the question. The benchmark therefore tests their application, rather than the discovery of physical laws.
Skateboarding provides recognizable examples whose names are easy to confuse. A kickflip and a heelflip differ in signed roll. A vertical spin differs from a roll even when both use the same numerical turn count. A late flip can share a rotation vector with a kickflip while requiring more airtime under the benchmark assumptions. These distinctions let one compact task expose several errors without requiring a visual scene.
This study contributes an explicit measurement account, controlled response experiments, and a simulated control-layer ablation. The controlled responses establish what the scoring instrument rewards. The simulated experiment shows how its dependencies affect a deliberately specified error population. Neither experiment measures the capability of a learned model. The scientific question is whether the score supports the interpretation that a reader might otherwise assign to it.
Our main finding is that partial correctness and complete consistency require separate reporting. Reversing signed rotations preserves considerable credit when arithmetic remains correct. Conversely, a small numeric error outside tolerance can remove boolean credit even when the booleans themselves match. This asymmetry is a design choice with practical consequences for comparing systems.
\section{Related work}
PHYRE evaluates physical reasoning through action in a simulated environment \citep{phyre}. Its puzzle setting connects interventions with observed outcomes. Skateboard instead provides a symbolic description and asks for a structured prediction. This removes perception and action selection from the measured construct. The narrower setting supports direct component checks, but it cannot establish competence at interaction with physical scenes.
CLEVRER studies descriptive, explanatory, predictive, and counterfactual reasoning over videos of object interactions \citep{clevrer}. Skateboard shares an interest in decomposing physical judgments. Its tasks do not require tracking objects or identifying collision causes. The distinction matters because success with supplied quantities and equations is weaker evidence than success with quantities inferred from a video.
BetterBench organizes benchmark quality around design, implementation, documentation, and maintenance \citep{betterbench}. The Agentic Benchmark Checklist emphasizes task validity, outcome validity, and transparent reporting \citep{abc}. We use these perspectives to examine the relation between the stated construct and the actual reward. A reference control establishes reachability. Trivial controls test whether non-solutions obtain credit. Deliberate errors test which components determine a score.
PayBench organizes evaluation around a construct, scenario coverage, controls, metrics, and caveats \citep{paybench}. We adopt that presentation logic. We do not claim its matched trap-and-lookalike design. The skateboard panel crosses trick definitions with gravity settings, while the response experiments hold tasks fixed and change answer components. These are controlled answer perturbations, not paired real-world safety scenarios.
The contribution is complementary to broader physical reasoning benchmarks. It gives a transparent account of a small symbolic instrument. It does not supply a new physics simulator, a model leaderboard, or evidence that its tasks predict success in sports. A stronger external claim would require independent tasks and observations beyond the supplied rules.
\section{Experiment design}
\subsection{The physical model}
The board and rider form a point mass during flight. Launch and landing occur at equal height. Gravity is constant and drag is absent. With upward speed $v$ and gravitational acceleration $g$, the model uses $t=2v/g$ and $h=v^2/(2g)$ for airtime and peak height. These expressions follow the standard constant-acceleration treatment \citep{openstax}. The benchmark does not estimate forces, foot contact, or angular acceleration.
The rotation vector is $r=(r_l,r_v,r_c)$ in full turns. Its components refer to the long, vertical, and cross axes. The view runs from tail toward nose. Kickflip direction is positive on the long axis; frontside is positive on the vertical axis; nose-forward motion is positive on the cross axis. Signs are conventions declared in the prompt, rather than facts inferred from camera position or rider stance.
The wheels-down label is true when long and cross turns are integers. Vertical turns do not change that label. Feasibility is true when exact airtime meets the trick's stated minimum. The two labels answer different questions. A trick can have enough time to complete its prescribed rotation while ending with an inverted board under the orientation rule. A correct answer must preserve that distinction.
This representation is intentionally coarse. General three-dimensional rotations depend on order. A vector of signed turn totals does not represent arbitrary rigid-body motion. The benchmark's integer rule is a defined endpoint convention, not a theorem about all composed rotations. The model also omits rider balance and the time needed to stabilize a landing. Its feasibility label concerns only the supplied duration threshold.
\subsection{Task construction and experimental unit}
The public panel contains \TaskCount{} tasks from \TrickCount{} trick definitions crossed with Earth, Moon, and Mars gravity settings. Every trick appears in every world. A seeded generator selects launch speeds from bands intended to include both sufficient and insufficient airtime. The public panel is fixed for this study. It is a census of this release, not a random sample of all skateboard situations.
The prompt states the trick description, equations, sign convention, required duration, and output schema. The analytic control extracts the numeric quantities and uses the public glossary to map a trick name to its rotation vector. It then computes both physics values and both booleans. This control has access to task semantics through a curated lookup. It does not demonstrate natural-language understanding, and it never reads a stored answer key when generating a response.
The design plan precedes the study's data files in version history. It is not an externally registered protocol. We inspected existing validation policies and a prior family report before writing it. Their results informed the controls, but none of their measurements enter this study's tables. A correction to the bootstrap implementation is recorded separately because its first implementation resampled within tricks rather than resampling whole tricks.
\subsection{Scripted controls and paired perturbations}
The analytic policy provides a positive control. Empty, refusal, and echo policies test obvious non-answers. The constant-zero policy supplies valid numeric and boolean types but ignores the question. The always-true policy preserves analytic rotations and physics while setting both booleans true. The reverse-sign policy preserves analytic physics and booleans while negating every rotation component. These policies isolate response properties, rather than approximating named models.
The sensitivity experiment multiplies both exact physics values by a fixed factor. It leaves analytic rotations and booleans unchanged. The factors span both sides of the relative tolerance boundary. This intervention is intentionally inconsistent with a single physical change to launch speed: time and height have different dependence on speed. It tests the scorer's numeric acceptance rule, not a new trajectory model.
All policies run through the public benchmark execution path with the same task panel and deterministic scorer. No external inference service is used. We repeat the entire condition grid and compare semantic records after excluding timing fields. Repetition checks execution stability; it does not create additional independent tasks or observations for uncertainty estimation.
\section{Coverage and task semantics}
\begin{figure}[ht]\centering
\begin{tikzpicture}\begin{axis}[ybar,ymin=0,ylabel={Task count},xtick={0,1,2},xticklabels={Earth,Moon,Mars},legend style={font=\footnotesize,at={(0.5,1.02)},anchor=south}]
\addplot table[x=x,y=n,col sep=comma]{figures/coverage.generated.csv};\addplot table[x=x,y=feasible,col sep=comma]{figures/coverage.generated.csv};\addplot table[x=x,y=wheels,col sep=comma]{figures/coverage.generated.csv};\legend{All,Feasible,Wheels down}
\end{axis}\end{tikzpicture}
\caption{Coverage counts describe the complete public panel. These are inventory totals, not estimated population rates. Equal world counts do not imply equal feasibility counts.}\label{fig:coverage}
\end{figure}

\begin{table}[ht]\centering\begin{tabular}{llll}\toprule
World & Tasks & Feasible count & Wheels-down count\\ \midrule
earth & 24 & 7 & 12\\
moon & 24 & 15 & 12\\
mars & 24 & 12 & 12\\
\bottomrule\end{tabular}
\caption{Complete public-panel inventory. Feasible and wheels-down counts can overlap.}\end{table}
Equal counts across worlds control one source of composition imbalance. They do not ensure equal difficulty or equal class balance. The launch-speed ranges differ by world, and the generator can favor underpowered launches where a trick's threshold permits them. Coverage must therefore be reported separately from outcomes. A lower score in one world would not, by itself, demonstrate weaker reasoning about that gravity value.
The glossary includes ordinary named tricks and incomplete-flip variants. These variants make orientation errors observable. A policy that always predicts wheels down can still solve the numeric part, but it cannot satisfy every orientation label. Likewise, the duration threshold distinguishes matching rotation vectors that have different timing assumptions. This supports separate tests for directional meaning, numeric substitution, and conditional judgment.
The coverage has important limits. Each trick-world combination appears with one public numeric instance. Alternative wording, rider stance, and camera transformations are absent. A solver can memorize glossary entries or the entire public panel. Fresh launch speeds would reduce exact-answer reuse while leaving names, templates, and rules familiar. Such recombination would not establish semantic generalization.
There is also a wording defect worth preserving in the validity account. Some incomplete-flip descriptions say the board comes down ``on its side.'' A half turn about its length instead inverts a flat board under the stated turn convention. The explicit integer rule still determines the scored boolean, but the prose can conflict with physical interpretation. This paper does not silently repair that wording or treat agreement with the key as independent validation of the description.
\section{Taxonomy of failure modes}
\paragraph{Axis and sign errors.} A solver can identify the amount of rotation while assigning the wrong sign or axis. The scorer checks the entire vector exactly. This makes the rotation component easy to audit. It also means a single sign error loses the same rotation credit as a completely incorrect vector. The reported score cannot reveal the size of a rotation error without the raw answer.
\paragraph{Numeric errors.} Airtime and height require different expressions. Confusing them, using the wrong gravity, or rounding too early can change both feasibility and component credit. The grader compares against unrounded physics when the task supplies valid speed and gravity metadata. Displayed reference answers are rounded, so they are not the final numeric authority for these comparisons.
\paragraph{Conditional errors.} The landing and duration labels use different evidence. A model can conflate sufficient time with wheels-down orientation. It can also apply a condition to its own incorrect intermediate values and produce an internally consistent but wrong answer. The simulated boolean-repair intervention directly tests the latter case.
\paragraph{Format errors.} The scorer requires an object with the specified fields, finite numbers, and boolean values. Missing fields or extra fields invalidate the answer. The extraction step can accept a surrounding fence and can locate JSON within text. Thus strict object shape does not imply strict enforcement of the prompt's request to return only JSON. These are separate properties of the response interface.
\paragraph{Score interpretation errors.} A high weighted average can conceal a low complete-answer fraction. Correct numbers can preserve most credit despite wrong rotations. Correct booleans can lose credit because a numeric prerequisite fails. Neither pattern is necessarily a grader defect. Both become misleading when users interpret the aggregate as the fraction of tasks solved.
\section{Metrics and scoring}
Let $R,T,H,L,F$ indicate correct rotations, airtime, height, orientation, and feasibility. Let $V$ indicate a valid response object. The task score is
\[S=V\{0.24R+0.20T+0.18H+0.19(TH)L+0.19(TH)F\}.\]
Here $L$ and $F$ compare booleans with reference labels. Their credit also requires both numeric checks. Incorrect rotations do not block either boolean term. The score therefore enforces a numeric prerequisite but not a rotation prerequisite. We retain these release rules rather than adjusting them to make the controls separate more sharply.
For numeric reference $x$, the acceptance condition is $|\hat{x}-x|\leq\max(0.05|x|,0.0005)+10^{-9}$. The absolute floor accommodates output rounded to three decimals. Consequently, a value can exceed the relative band and still pass when its magnitude is small. This is not an exclusion rule. All such cases remain in the sensitivity analysis.
The primary outcome is the mean of $S$. Complete-answer accuracy is the fraction with $S=1$. Component credits accompany these measures in the released data. A component credit is a weighted contribution, not a normalized component accuracy. The mean can rise through improvements that never produce a complete answer, so it should not replace the secondary outcome.
\subsection{Uncertainty and pairing}
Fixed-policy mean intervals use a percentile bootstrap over whole tricks with \BootstrapDraws{} draws and nominal 95\% coverage. Each sampled trick carries all its world instances. This retains dependence among tasks sharing a rotation definition. Complete-answer fractions include Wilson intervals over tasks for comparison, plus trick-cluster intervals in the data. Wilson intervals assume independent Bernoulli trials and therefore do not account for shared trick structure. They are conditional descriptive summaries, not evidence of population representativeness.
The simulated experiment instead resamples agent seeds. Each simulated agent retains its complete panel and its paired control outcomes. Control gains are calculated within agent before resampling. This estimates variability under the declared simulated response process, conditional on the fixed public tasks. It says nothing about the distribution of real models or users.
Boundary bootstrap intervals can collapse when every retained cluster has the same outcome. That behavior does not prove absence of uncertainty outside the panel. We report the Wilson counterpart for fixed-policy full-credit outcomes and state the conditional scope of all intervals. We perform no significance test, model ranking, or post hoc task selection.
\section{Results}
\subsection{Analytic ceiling and trivial controls}
\begin{table}[ht]\centering\small\begin{tabular}{lll}\toprule
Policy & Mean [cluster CI] & Full [Wilson CI]\\ \midrule
analytic & 1.000 [1.000, 1.000] & 1.000 [0.949, 1.000]\\
empty & 0.000 [0.000, 0.000] & 0.000 [0.000, 0.051]\\
refuse & 0.000 [0.000, 0.000] & 0.000 [0.000, 0.051]\\
echo & 0.000 [0.000, 0.000] & 0.000 [0.000, 0.051]\\
constant-zero & 0.020 [0.000, 0.050] & 0.000 [0.000, 0.051]\\
always-true & 0.805 [0.760, 0.847] & 0.250 [0.164, 0.361]\\
reverse-sign & 0.780 [0.760, 0.810] & 0.083 [0.039, 0.170]\\
\bottomrule\end{tabular}
\caption{Scripted controls. Brackets are 95\% intervals. Means resample tricks; full-credit intervals use Wilson's task-level model.}\end{table}
\begin{figure}[ht]\centering
\begin{tikzpicture}\begin{axis}[ymin=0,ymax=1.1,ylabel={Mean score},xtick={0,1,2,3,4,5,6},xticklabels={Analytic,Empty,Refuse,Echo,Constant,True,Reverse},x tick label style={rotate=25,anchor=east}]
\addplot+[color=okblue,mark=*,error bars/.cd,y dir=both,y explicit] table[x=x,y=y,y error minus=minus,y error plus=plus,col sep=comma]{figures/baselines.generated.csv};
\end{axis}\end{tikzpicture}
\caption{Scripted controls separate full solutions from partial credit. Bars show 95\% trick-cluster bootstrap intervals. The sign-reversed control preserves correct arithmetic.}\label{fig:baselines}
\end{figure}

The analytic control reaches mean \OracleMean{} and full-credit fraction \OracleFull{}. This supports the feasibility of the answer format and the supplied mathematical rules on the public panel. It is weaker than independent validation: the control uses the same public glossary that underlies task construction. A shared glossary error could survive this check.
Empty, refusal, and echo outputs receive no credit. Their complete-answer uncertainty appears in the table. The constant-zero control obtains mean \ConstantMean{}, despite never producing a complete answer. That partial credit comes from tasks whose rotation vectors are zero. The result does not contradict the planned hypothesis about complete answers, but it rules out describing every task-ignoring policy as a zero-score baseline.
The always-true control obtains mean \AlwaysMean{} with full-credit fraction \AlwaysFull{}. It knows the rotations and physics by construction, so it is not a blind affirmative baseline. Its residual failure measures the need to distinguish orientation and feasibility labels. Reporting this policy as a learned reasoner would misstate the experiment.
The reverse-sign control obtains mean \ReverseMean{}, while its full-credit fraction is \ReverseFull{}. Zero vectors are unchanged by negation. Other rotations fail while numeric and boolean terms remain correct. This demonstrates a measurement asymmetry: a severe directional error can coexist with a large aggregate score. It supports reporting full answers and component errors alongside the headline mean.
\subsection{Numeric sensitivity and the rounding exception}
\begin{table}[ht]\centering\small\begin{tabular}{lll}\toprule
Physics factor & Mean [cluster CI] & Full [Wilson CI]\\ \midrule
0.94 & 0.243 [0.240, 0.248] & 0.000 [0.000, 0.051]\\
0.96 & 1.000 [1.000, 1.000] & 1.000 [0.949, 1.000]\\
1 & 1.000 [1.000, 1.000] & 1.000 [0.949, 1.000]\\
1.04 & 1.000 [1.000, 1.000] & 1.000 [0.949, 1.000]\\
1.06 & 0.243 [0.240, 0.248] & 0.000 [0.000, 0.051]\\
1.1 & 0.240 [0.240, 0.240] & 0.000 [0.000, 0.051]\\
\bottomrule\end{tabular}
\caption{Controlled numeric perturbations. The multiplier changes both exact physics values. Brackets are 95\% intervals with the same units as Table 2.}\end{table}
\begin{figure}[ht]\centering
\begin{tikzpicture}\begin{axis}[ymin=0,ymax=1.1,xlabel={Multiplier on both exact physics values},ylabel={Mean score}]
\addplot+[color=okorange,mark=*,error bars/.cd,y dir=both,y explicit] table[x=x,y=y,y error minus=minus,y error plus=plus,col sep=comma]{figures/sensitivity.generated.csv};
\end{axis}\end{tikzpicture}
\caption{Physics tolerance creates a sharp score boundary. Bars show 95\% trick-cluster bootstrap intervals. The absolute rounding floor preserves a small amount of height credit outside the relative band.}\label{fig:sensitivity}
\end{figure}

The sensitivity curve changes sharply across the relative tolerance band. Outside it, boolean credit disappears because the numeric prerequisite fails. However, the planned prediction that all physics credit disappears is too strong. The upper outside condition retains mean \SensitivityOutside{}. A small reference height remains within the absolute rounding floor, preserving height credit even though its relative error exceeds the nominal band.
This exception is useful evidence about the instrument. A claim that the grader uses only relative error would be false. The absolute floor has a legitimate purpose, but it changes the effective tolerance for small values. Future studies should report numeric magnitude distributions and not infer an identical acceptance region from a single percentage tolerance.
The experiment does not measure model sensitivity to a perturbed prompt. It modifies answers after analytic computation. The causal interpretation is restricted to the mapping from submitted values to scores. A learned solver might change several fields together or produce malformed output, neither of which this numeric-only intervention models.
\section{Extra virtual experiment: simulated control ablation}
\subsection{Simulated population and intervention}
The simulated population contains \SimAgentCount{} seeded agents at each probability and control setting. Each simulated agent starts with an analytic answer. Independent seeded draws decide whether to corrupt its rotation vector, airtime, and height. Rotation corruption adds a half turn to the long component. Numeric corruption multiplies the selected rounded value by \SimMultiplier{}. The same draws are reused across corruption probabilities and control layers. This provides paired comparisons within each simulated agent and task.
The simulated no-repair condition leaves these errors in place. Simulated physics repair recomputes both numeric values from the prompt but leaves rotations and booleans unchanged. Simulated boolean repair recomputes orientation from the submitted rotations and feasibility from submitted airtime. It does not correct either source quantity. This separates recovery of correct values from enforcement of internal consistency.
The starting booleans are analytic in every simulated condition. Boolean repair therefore has no initially wrong boolean to rescue. This deliberately isolates error propagation, rather than estimating the usefulness of a general repair system. A population with independent boolean mistakes could yield a different conclusion. The intervention is transparent enough to make that dependency explicit.
\begin{figure}[ht]\centering
\begin{tikzpicture}\begin{axis}[ymin=0,ymax=1.1,xlabel={Simulated corruption probability},ylabel={Simulated mean score},legend style={font=\footnotesize,at={(0.02,0.02)},anchor=south west}]
\addplot+[color=black,mark=*,error bars/.cd,y dir=both,y explicit] table[x=x,y=y,y error minus=minus,y error plus=plus,col sep=comma]{figures/sim-none.generated.csv};
\addplot+[color=okblue,mark=*,error bars/.cd,y dir=both,y explicit] table[x=x,y=y,y error minus=minus,y error plus=plus,col sep=comma]{figures/sim-physics.generated.csv};
\addplot+[color=okorange,mark=*,error bars/.cd,y dir=both,y explicit] table[x=x,y=y,y error minus=minus,y error plus=plus,col sep=comma]{figures/sim-boolean.generated.csv};
\legend{No repair,Physics repair,Boolean repair}
\end{axis}\end{tikzpicture}
\caption{Simulated population sweep with paired random draws. Bars show 95\% bootstrap intervals over agent seeds. Simulated physics repair restores arithmetic; simulated boolean repair can propagate rotation errors.}\label{fig:population}
\end{figure}

\begin{table}[ht]\centering\footnotesize\begin{tabular}{llll}\toprule
Corruption & Simulated control & Mean [agent CI] & Full [agent CI]\\ \midrule
0 & none & 1.000 [1.000, 1.000] & 1.000 [1.000, 1.000]\\
0 & physics & 1.000 [1.000, 1.000] & 1.000 [1.000, 1.000]\\
0 & boolean & 1.000 [1.000, 1.000] & 1.000 [1.000, 1.000]\\
0.25 & none & 0.674 [0.663, 0.686] & 0.417 [0.398, 0.438]\\
0.25 & physics & 0.941 [0.934, 0.948] & 0.755 [0.726, 0.783]\\
0.25 & boolean & 0.651 [0.640, 0.663] & 0.417 [0.398, 0.438]\\
0.5 & none & 0.401 [0.384, 0.416] & 0.116 [0.101, 0.132]\\
0.5 & physics & 0.880 [0.873, 0.887] & 0.501 [0.473, 0.531]\\
0.5 & boolean & 0.379 [0.364, 0.393] & 0.116 [0.101, 0.132]\\
0.75 & none & 0.175 [0.163, 0.186] & 0.017 [0.010, 0.024]\\
0.75 & physics & 0.822 [0.818, 0.826] & 0.259 [0.243, 0.274]\\
0.75 & boolean & 0.168 [0.157, 0.178] & 0.017 [0.010, 0.024]\\
1 & none & 0.000 [0.000, 0.000] & 0.000 [0.000, 0.000]\\
1 & physics & 0.760 [0.760, 0.760] & 0.000 [0.000, 0.000]\\
1 & boolean & 0.000 [0.000, 0.000] & 0.000 [0.000, 0.000]\\
\bottomrule\end{tabular}
\caption{Simulated population outcomes. Every bracket is a 95\% bootstrap interval over simulated agent seeds. Full means complete-answer fraction.}\end{table}
\subsection{Simulated outcomes and paired effects}
At the middle corruption setting, the simulated no-repair mean is \Simnone{}. Simulated physics repair reaches \Simphysics{}, with paired gain \Gainphysics{}. Simulated boolean repair reaches \Simboolean{}, with paired gain \Gainboolean{}. These intervals describe the simulated agents, conditional on the task panel. They must not be read as model performance intervals.
Simulated physics repair restores both arithmetic terms and their boolean prerequisites. Its gain therefore exceeds the direct numeric credit that it changes. Simulated boolean repair can flip a correct orientation label to match a corrupted rotation. The result can become more internally consistent while moving farther from the reference answer. The scorer rewards reference agreement, not consistency with every submitted intermediate result.
The simulated full-credit fractions for boolean repair and no repair coincide in this design. A response with corrupted rotations already fails the full-answer condition. Changing its orientation cannot restore full credit. A response with corrupted physics also fails that condition. This null on complete answers coexists with a negative effect on the weighted mean. The two outcomes answer different questions.
\begin{table}[ht]\centering\small\begin{tabular}{lll}\toprule
Corruption & Simulated control & Paired gain [agent CI]\\ \midrule
0 & physics & 0.000 [0.000, 0.000]\\
0 & boolean & 0.000 [0.000, 0.000]\\
0.25 & physics & 0.267 [0.254, 0.279]\\
0.25 & boolean & -0.023 [-0.027, -0.019]\\
0.5 & physics & 0.480 [0.463, 0.498]\\
0.5 & boolean & -0.021 [-0.025, -0.018]\\
0.75 & physics & 0.647 [0.636, 0.658]\\
0.75 & boolean & -0.008 [-0.009, -0.006]\\
1 & physics & 0.760 [0.760, 0.760]\\
1 & boolean & 0.000 [0.000, 0.000]\\
\bottomrule\end{tabular}
\caption{Paired simulated control gains relative to no repair. Intervals resample complete simulated agents and retain condition pairing.}\end{table}
The simulated curves are expected from the declared mechanism, rather than discoveries about natural populations. Their value is diagnostic: they show how a control can improve one dependency, propagate another error, or leave complete-answer accuracy unchanged. They also identify which assumptions a later learned-model study must test before transferring these conclusions.
\section{Roadmap}
The first priority is independent task validation. Reviewers should resolve conflicts between prose, rotation conventions, and endpoint rules before using the benchmark for external claims. A separate expert audit should distinguish stylized definitions from descriptions that imply real board motion. Corrections should create a versioned evaluation set, because changes to wording can change what a solver is asked to do.
The next priority is a held-out semantic design. New numbers alone preserve the same templates and lookup structure. Useful extensions include paraphrased rotation descriptions, unfamiliar compositions with explicit order, and separate tests of inference from diagrams. Each extension changes the construct and needs its own positive controls. They should not be pooled into a single score without a clear weighting argument.
A learned-model study should compare complete answers, component errors, and invalid outputs under a fixed prompt. It should also compare repair interventions on naturally occurring mistakes. The present simulated population holds booleans correct initially; a later study must measure their actual dependence on intermediate numeric and rotation errors. Free hosted inference was not configured here, and provider availability was not probed.
Finally, calibration near the decision boundary needs new tasks designed for that purpose. The current speed bands leave margins around many feasibility thresholds. Correct results on this panel cannot establish reliable behavior for almost-feasible trajectories. A future boundary panel should retain exact numeric references and disclose how rounding affects each case.
\section{Caveats and scope of evidence}
The observations support claims about a fixed scoring instrument and scripted response families. They do not show that current language models reason well about skateboarding. No human baseline, motion capture, robot experiment, or field outcome is present. No learned-model score is published. We do not claim that the proposed metric predicts coaching quality or safe physical execution.
The task descriptions carry assumptions about timing and orientation. Those assumptions are not fitted to skaters. In particular, a minimum duration is a benchmark constraint rather than a measured biomechanical limit. The inclusion of lunar and Martian settings expands symbolic gravity variation without supplying evidence about actual skateboarding in those environments.
The public panel has shared templates and glossary entries. The cluster analysis retains dependence within tricks, but it cannot remove all dependence created by common wording or the same equations. Confidence intervals quantify uncertainty under stated resampling models. They do not compensate for limited coverage or an incorrect construct definition.
The simulated response process also has narrow scope. Corruptions are independent by field and task under seeded hashing. Real model errors can cluster by vocabulary, arithmetic difficulty, or reasoning strategy. The repair layers have privileged access to explicit quantities and a deterministic rule. Their measured effects should be transferred only after checking those assumptions.
The whole grid contains \ConditionCount{} conditions and \ObservationCount{} scored responses per pass. The repeated pass agrees in semantic hashes. These counts represent computation, not a larger number of independent task families. Repeated deterministic agreement supports reproducibility but cannot detect a shared mistake in the glossary, analytic policy, and scorer.
\section{Availability, citation, and authors}
Code, data, the committed design plan, and reproduction instructions are available at \url{https://github.com/pooriaarab/benchmark-skateboard}. The paper and generated figures use CC-BY-4.0. The software retains its existing license. Cite this paper by its title, author, and versioned source record; no new publication DOI is claimed.
Pooria Arab is an independent researcher. AI assistance supported implementation and drafting. The author-defined study uses no paid inference or human participants. The packet is a research draft for review, not a venue submission or an independently certified benchmark.
\bibliographystyle{plainnat}
\bibliography{references}
\end{document}
